\section{Instruction Set and Core Structures}
\label{sec:instruction_set}

\subsection{The IsalSR Alphabet $\Sigma_{\mathrm{SR}}$}

IsalSR uses a \textbf{two-tier token encoding}: single-character tokens for
movement and edge operations, and two-character compound tokens for labeled
node insertion.

\begin{definition}[IsalSR Alphabet]
The alphabet $\Sigma_{\mathrm{SR}}$ consists of:
\begin{itemize}[nosep]
  \item \textbf{Movement} (single-char): $\{N, P, n, p\}$ --- move primary/secondary
    pointer next/prev in the CDLL.
  \item \textbf{Edge creation} (single-char): $\{C, c\}$ --- directed edge
    primary$\to$secondary / secondary$\to$primary, with DAG cycle check.
  \item \textbf{No-op} (single-char): $\{W\}$.
  \item \textbf{Node insertion} (two-char): $V\ell$ or $v\ell$ where $\ell$ is a
    label character from
    $\mathcal{L} = \{+, *, -, /, s, c, e, l, r, \hat{}, a, k\}$.
\end{itemize}
Total vocabulary: 7 single-char $+$ 24 compound $= 31$ tokens (configurable
subset per experiment).
\end{definition}

\begin{table}[h]
\centering
\caption{Operation types and their label characters.}
\label{tab:operations}
\begin{tabular}{llcl}
\toprule
Label & NodeType & Arity & Description \\
\midrule
$+$ & ADD & Variable ($\geq 2$) & Sum of all inputs \\
$*$ & MUL & Variable ($\geq 2$) & Product of all inputs \\
$-$ & SUB & Binary (2) & Subtraction \\
$/$ & DIV & Binary (2) & Protected division \\
$s$ & SIN & Unary (1) & Sine \\
$c$ & COS & Unary (1) & Cosine \\
$e$ & EXP & Unary (1) & Protected exponential \\
$l$ & LOG & Unary (1) & Protected logarithm \\
$r$ & SQRT & Unary (1) & Protected square root \\
$\hat{}$ & POW & Binary (2) & Protected power \\
$a$ & ABS & Unary (1) & Absolute value \\
$k$ & CONST & Leaf (0) & Learnable constant \\
\bottomrule
\end{tabular}
\end{table}

\subsection{Initial State}

For $m$ input variables $x_1, \ldots, x_m$:
\begin{itemize}[nosep]
  \item The DAG contains $m$ VAR nodes (IDs $0, 1, \ldots, m-1$), no edges.
  \item The CDLL contains $m$ entries $[x_1, x_2, \ldots, x_m]$ in order.
  \item Both pointers (primary, secondary) start on the CDLL node for $x_1$.
\end{itemize}

Variables are \textbf{pre-numbered and distinguishable}: $x_1 \neq x_2$ under
isomorphism. This eliminates isomorphism ambiguity over variable assignments
and is the key simplification that allows the canonical string to be computed
from $x_1$ only.

\subsection{Edge Semantics and the DAG Constraint}

An edge $u \to v$ means ``$u$ provides input to $v$'' (data flow direction).
For $\sin(x)$: edge $x \to \sin$. For $x + y$: edges $x \to +$ and $y \to +$.

The $C$ instruction creates a directed edge from the primary pointer's graph
node to the secondary pointer's graph node. Before adding any edge via $C/c$,
the system checks whether the edge would create a cycle (via BFS reachability).
If a cycle would result, the instruction is \textbf{silently skipped}.
$V/v$ insertions never create cycles because new nodes have no outgoing edges.

\subsection{Core Data Structures}

\textbf{LabeledDAG}: Directed acyclic graph with dual adjacency lists
(in-edges and out-edges), node labels ($\texttt{NodeType}$ enum), and per-node
metadata (variable index, constant value). Supports $O(V+E)$ cycle detection,
Kahn's topological sort, and label-aware backtracking isomorphism.

\textbf{CircularDoublyLinkedList (CDLL)}: Array-backed circular list with $O(1)$
insert-after, remove, and traversal. Reused verbatim from IsalGraph
\cite{lopezrubio2025}. Pointers are CDLL node indices (not graph node IDs);
conversion via \texttt{get\_value(ptr)}.
